\documentclass[./Main.tex]{subfiles}
\graphicspath{{./Fig/}}
\begin{document}

\setcounter{chapter}{-1}
\chapter{このテンプレートの説明}
\section{修論の構成}
修論，卒論ともに通常の学術論文と同じように

\begin{enumerate}
    \item 序論（背景，既存の研究のレビューと目的）
    \item 研究手法
    \item 結果と考察
    \item 結論
    \item その他（参考文献，謝辞，付録等。ただし，参考文献は必須）
\end{enumerate}
の順序で書きます。（実際には「その他」という章は普通作りません）それに対応して，このテンプレートではTeXのソースファイルを
\begin{itemize}
    \item introduction.tex
    \item solution.tex
    \item result.tex
    \item conclusion.tex
\end{itemize}
にしています。参考文献は最後にまとめて出力するようにしています。

TeXの場合，ソースファイルをどう分割しても特に問題ないので，手法や結果がながければ，さらに細かく分割したほうが見やすいです。

\subsection{LaTeXのエンジン}
このテンプレートはplatexでコンパイル，dvipdfmxでPDF化するようにしています。

\subsection{ファイル分割の詳細}
通常，ソースファイルを分割する場合，\verb|\include \input|命令を使います。ですが，ここではsubfilesパッケージを使っています。subfilesパッケージを使うと，分割した各ソースファイルを単独でコンパイルできるようになります。例えば，introduction.texを単独でコンパイルすると，introduction.texだけで完結したPDFが生成されます。これにより，部分的にコンパイルして確認したい場合に便利です。

subfilesでファイルを分割する場合，Mainのソースファイルでは，\verb|\subfile{ファイル名.tex}|とします。\verb|\input|等と異なり，拡張子の.texまで必要です。分割したファイル側には，\verb|\documentclass[./Main.tex]{subfiles}|をはじめに書きます。Main.texが親ファイルであることを指定するためです。このテンプレートのファイルを確認してください。分割したファイル単独でコンパイルできるようになります。コード\ref{code_subfile}に分割したTeXのソース例を示します。最後のifなんちゃらは，大本のMain.texをコンパイルしたときには無視されます。

\begin{lstlisting}[caption=subfileで分割したTeXのソース, label=code_subfile]
\documentclass[./Main.tex]{subfiles}
\graphicspath{{./Fig/}}
\begin{document}

% この間に本文を書く

\ifSubfilesClassLoaded{
\printbibliography[title = 参考文献]
}{}

\end{document}
\end{lstlisting}

\subsection{参考文献の出力}
参考文献はbibtexで出力しています。preamble.texを確認してもらえれば分かります。
自分でコンパイルするときには以下の流れになります。MainのファイルがMain.texとして，

\begin{verbatim}
platex Main.tex
bibtex Main
platex Main.tex
platex Main.tex
dvipdfmx Main.dvi
\end{verbatim}
の順でコンパイルします。最初のplatexでauxファイルが生成され，bibtexで参考文献データが処理され，その後のplatexで参考文献が出力されます。2回platexを回すのは，参照番号などを正しく反映させるためです。


\section{LaTeXの各パッケージと使い方}

\subsection{zebra-goodies}
todoやcommentを本文に挿入するためのパッケージ。
TeXのソースに
\begin{verbatim}
\todo{ここにTODOを書く}
\comment{XXについて要確認}
\end{verbatim}
という感じで，コマンドを入れると，PDFに赤字でTODOやコメントが挿入されます。

原稿が完成して，todoやコメントを一括で消したいときには
\begin{verbatim}
    \usepackage[final]{zebra-goodies}
\end{verbatim}
でコンパイルすれば，todoやcommentが消えます。

\subsection{cleverref}
図番号・表番号を参照するのに通常は\verb|\ref{}|を使いますが，\verb|\cref{}|を使うと，自動的に「図」「表」などの接頭語が付きます。例えば，\verb|\cref{fig:example}|とすると，「図1」のように出力されます。



\ifSubfilesClassLoaded{
\printbibliography[title=参考文献]
}{}

\end{document}